\section{A near-optimal bounded classical formula}
\label{sec:classical-optimization}

We now construct a bounded classical output whose rooted energy is
arbitrarily close to $P_*$. The construction has two parts. A scalar
Parisi diffusion first gives a finite function of independent Gaussians
with an explicitly computable value. We realize those Gaussians as
positions in the recursive space of
Section~\ref{subsec:classical-gaussian-sector}. Bounded approximation
and type exclusions then put the formula in the class required for
quantum synthesis. These last two operations will also serve the
independent positive-temperature construction in
Appendix~\ref{sec:positive-temperature-alternative}.

\subsection{The scalar input}

We use the following results from the complete companion
\emph{Full support of the zero-temperature Sherrington--Kirkpatrick
order parameter}~\cite{SKSupport2026}. We state the normalization and
the finite formula explicitly because the rooted realization below
depends on their exact form.

Let $\mathcal U$ consist of the nonnegative, nondecreasing,
right-continuous functions $\gamma:[0,1)\to[0,\infty)$ with
$\int_0^1\gamma(t)\,\dd t<\infty$. Define the Parisi value function
by the controlled heat equation
\[
 \Phi_t=-\tfrac12(\Phi_{xx}+\gamma(t)\Phi_x^2),
 \qquad \Phi(1,x)=|x|,
\]
in its integrable-coefficient control interpretation, and put
\[
 \mathcal P(\gamma)=\Phi_\gamma(0,0)
                 -\tfrac12\int_0^1t\gamma(t)\,\dd t.
\]
The ground-state formula and attainment of Auffinger and
Chen~\cite[Theorem~1 and the proof of Lemma~3]{AC2017} identify
$\min_{\gamma\in\mathcal U}\mathcal P(\gamma)$ with the constant
$P_*$ in~\eqref{eq:parisi-constant}. Fix any such minimizer.
Throughout this section $\gamma(t)$ is this scalar coefficient,
not a vector of cost angles.

For clarity, the covariance of the unordered-pair Hamiltonian is
$\E h_{n,J}(\sigma^1)h_{n,J}(\sigma^2)=nR_{12}^2/2-1/2$, where
$R_{12}=n^{-1}\sum_i\sigma_i^1\sigma_i^2$. Adding a centered
Gaussian of variance $1/2$, constant across configurations, gives
covariance $n\xi(R_{12})$ without changing the expected maximum.
This is the normalization of the cited formula.

Write $u=\Phi_{\gamma,x}$, $a=\Phi_{\gamma,xx}$ and let
\[
 \dd X_t=\gamma(t)u(t,X_t)\,\dd t+\dd B_t,\qquad X_0=0.
\]
The companion proves full relative Stieltjes support on $[0,1)$,
with the convention $\mu_\gamma([0,t])=\gamma(t)$, and gives
\begin{equation}
 \E u(t,X_t)^2=t,\qquad
 \E a(t,X_t)^2=1,\qquad
 P_*=\int_0^1\E a(t,X_t)\,\dd t.
 \label{eq:scalar-consistency-value}
\end{equation}
Here $u$ is odd, $a$ is even and nonnegative, and $|u|\le1$.
All their spatial derivatives are bounded on every fixed strip
ending before one. These are
Theorem~1.1, Lemmas~2.1--2.2 and Corollary~7.1 of
\cite{SKSupport2026}; their complete analytic proofs belong to that
companion, not to the rooted-space argument below. In particular,
Cauchy--Schwarz gives
\begin{equation}
 0\le P_*-\int_0^T\E a(t,X_t)\,\dd t\le1-T,
 \qquad 0<T<1.
 \label{eq:scalar-terminal-loss}
\end{equation}
No smooth density of $\mu_\gamma$ or bound uniform as $T\uparrow1$
is used here.

The use of the Parisi gradient and a normalized Euler martingale
follows the classical constructions of
Montanari~\cite[Section~3.1]{Montanari2019} and El Alaoui,
Montanari, and Sellke~\cite[Proposition~5.3]{AMS2020}.
The precise shifted-coefficient identity needed here is the scalar
companion input stated below.

For fixed $0<T<1$ and an integer $N\ge1$, set $\Delta=T/N$ and
$t_j=j\Delta$. On independent standard normal variables
$Z_1,\ldots,Z_{N+1}$, define
\begin{align}
 Y_0&=0,\qquad
 Y_j=Y_{j-1}+\sqrt\Delta\,Z_j
       +\Delta\gamma(t_{j-1})u(t_{j-1},Y_{j-1}), \label{eq:scalar-euler}\\
 c_j&=\bigl(\E a(t_{j-1},Y_{j-1})^2\bigr)^{-1/2},
 \qquad
 g_j=c_j a(t_{j-1},Y_{j-1})Z_j,\qquad 1\le j\le N.
 \label{eq:scalar-columns}
\end{align}
For all sufficiently large $N$, each $c_j$ is finite and positive.
With $\Phi_{\rm G}$ the standard normal distribution function, put
\begin{align}
 H_N&=u(T,Y_N)+2\Phi_{\rm G}(Z_{N+1})-1,
 &F_N&=\operatorname{sgn}(H_N), \label{eq:scalar-rounding}\\
 A_j&=\E[F_Ng_j]\quad(1\le j\le N),&
 B_i&=\E[F_NZ_i]\quad(1\le i\le N+1),\qquad A_0=0 .
 \label{eq:scalar-coefficients}
\end{align}
Conditional on $Z_1,\ldots,Z_N$, the added variable
$2\Phi_{\rm G}(Z_{N+1})-1$ is uniform on $(-1,1)$. Since $|u|\le1$,
\[
 \E[F_N\mid Z_1,\ldots,Z_N]=u(T,Y_N).
\]
Thus sign rounding preserves the terminal gradient as its conditional
mean. The same uniformity makes the value of the sign at zero
immaterial. Proposition~7.2 of~\cite{SKSupport2026} states that
\begin{equation}
 \lim_{N\to\infty}\sum_{j=0}^N A_jB_{j+1}
       =\int_0^T\E a(t,X_t)\,\dd t.
 \label{eq:scalar-coefficient-limit}
\end{equation}
The mesh limit in this assertion is taken at fixed $T$; only
afterwards is $T$ allowed to approach one. This finite scalar
identity, rather than a finite-degree message-passing theorem, is
the input to the next construction.

\subsection{Exact realization in the rooted Gaussian space}

Recall that a real centered child vector $f$ defines a Gaussian
position $\mathsf G(f)=q(f)$ at its parent, with covariance
$\E[\mathsf G(f)\mathsf G(g)]=\braket{f}{g}$.
The parent seed is independent of these child-Fock coordinates.
These two facts allow the scalar variables above to be constructed
successively at finite height.

\begin{lemma}[A finite Gaussian formula and its rooted energy]
\label{lem:gaussian-rooted-realization}
Fix $T,N$ and the scalar data in
\eqref{eq:scalar-euler}--\eqref{eq:scalar-coefficients}, with all
$c_j$ finite and positive. There is a real odd finite descendant
formula $F_N$ of height at most $N+1$ in the classical rooted space
whose Gaussian coordinates have exactly the stated joint law and
whose energy satisfies
\begin{equation}
 \cE(F_N)=\sum_{j=0}^N A_j B_{j+1}.
 \label{eq:scalar-rooted-energy}
\end{equation}
The formula is bounded by one, although its transmitted columns
need not be bounded.
\end{lemma}

\begin{proof}
Let $S$ be a standard Gaussian in the local root-seed factor and
start with $g_0=S$. Recursively set
\[
 Z_{j+1}=\mathsf G(g_j),\qquad 0\le j\le N,
\]
and use~\eqref{eq:scalar-euler}--\eqref{eq:scalar-columns} to form
$Y_j,g_j$ as soon as $Z_j$ is available. The constants $c_j$ here
are those defined under the independent Gaussian law above.
We will verify that the recursively constructed variables have
precisely that law.

At each step identify all earlier vectors by the compatible
empty-descendant embeddings of Section~\ref{sec:foundations}.
These embeddings preserve the earlier Gaussian positions:
lifting an embedded child vector gives the embedded copy of its
previous lift. Thus the recursion uses the same earlier variables
at the next height. Inductively $g_j$ has height at most $j$ and
$Z_{j+1}$ has height at most $j+1$. All vectors can be compared
in the same finite-height space whenever an inner product is taken.

We show inductively that $g_0,\ldots,g_j$ are real centered
orthonormal vectors. The claim holds for $g_0$. If it holds through
$g_{j-1}$, the positions $Z_1,\ldots,Z_j$ are jointly centered
Gaussian with identity covariance, hence independent standard
normals. For $j\ge1$, the coefficient
$c_j a(t_{j-1},Y_{j-1})$ is measurable in
$Z_1,\ldots,Z_{j-1}$. Conditional mean zero of $Z_j$ proves that
$g_j$ is centered and orthogonal to all earlier $g_i$ with $i\ge1$.
The normalization in~\eqref{eq:scalar-columns} gives
$\norm{g_j}=1$. The receiving-root seed $S$ belongs to a different
tensor factor from every child-Fock position, so it is independent
of all these positions. In particular
$\braket{g_0}{g_j}=0$.
The child directions themselves may contain seeds at their own
roots; this does not change the independence at the receiving root.
This closes the induction and also proves that
$Z_1,\ldots,Z_{N+1}$ have the required joint law. The coefficients
are bounded at the fixed grid times, so every constructed column
has moments of every finite order.

The map $u(t,\cdot)$ is odd and $a(t,\cdot)$ is even. Consequently
each $Y_j$ and $g_j$ is odd under simultaneous seed reversal.
Since $\Phi_{\rm G}(-z)=1-\Phi_{\rm G}(z)$, the same is true of
$F_N$ outside its null zero set. Hence $F_N$ is centered.

It remains to compute its energy, not just its distribution.
The vector $F_N$ is constant in the root-seed factor and depends
only on the finite orthonormal family of Gaussian modes
$\mathsf G(g_0),\ldots,\mathsf G(g_N)$.
The adjoint root shift extracts precisely its constant-root-seed
one-particle component. Since $q(g)=a^*(g)+a(g)$ has variance
$\norm g^2$, this component is
\[
 \cL^*F_N=\sum_{j=0}^N\E[F_NZ_{j+1}]\,g_j
        =\sum_{j=0}^N B_{j+1}g_j.
\]
There are no other first-chaos coefficients: $F_N$ is independent
of Gaussian modes orthogonal to the displayed family.
Moreover $F_N$ is independent of the receiving-root seed, so
$\braket{F_N}{g_0}=0=A_0$. Inner products with the other $g_j$
are exactly the scalar coefficients $A_j$. All vectors are real,
and therefore
\[
 \cE(F_N)=\re\braket{F_N}{\cL F_N}
        =\re\braket{\cL^*F_N}{F_N}
        =\sum_{j=0}^N A_jB_{j+1}.
\]
This proves the claimed identity, with the Gaussian variance and
root-shift normalizations used throughout the paper.
\end{proof}

Combining~\eqref{eq:scalar-terminal-loss},
\eqref{eq:scalar-coefficient-limit} and
\eqref{eq:scalar-rooted-energy} produces a bounded near-optimal
rooted vector by first choosing $T$ and then a finite $N$.
We next make its output smooth. At these fixed parameters set
$F_{N,k}=\tanh(kH_N)$. The null-set observation after
\eqref{eq:scalar-coefficients} and dominated convergence give
$F_{N,k}\to F_N$ in $L^2$. The root-shift continuity estimate is
\begin{equation}
 |\cE(f)-\cE(g)|
 \le(\norm f+\norm g)\norm{f-g}.
 \label{eq:energy-l2-continuity}
\end{equation}
Both outputs have norm at most one. Thus for every
$\varepsilon>0$ we may choose finite $T,N,k$ so that
\begin{equation}
 z=F_{N,k},\qquad |z|\le1,\qquad
 \cE(z)\ge P_*-\varepsilon.
 \label{eq:scalar-near-optimal-smooth}
\end{equation}

For the already fixed $T,N,k$, every transmitted-column map and
the final output map is jointly odd and smooth, with all
derivatives of polynomial growth in its finite list of odd
inputs. Indeed the Euler maps use finitely many fixed-time
bounded spatial derivatives of $u,a$, followed in a column by
multiplication by one Gaussian coordinate. Local intermediate
operations may be composed into those maps. The final sign
smoothing has bounded derivatives of every order at fixed $k$.

Two further operations are still necessary: bounding every
transmitted column and imposing type exclusions. We perform
them directly in the limiting Hilbert space. No finite-degree
edge-energy convergence is asserted for the discontinuous raw
$F_N$.

\subsection{Bounded trigonometric approximation}

For use with either classical construction, let
$\pi(x)=\max(-1,\min(x,1))$. The next lemma applies to any finite
descendant formula in its stated class. In the present application
$\pi(z)=z$ by~\eqref{eq:scalar-near-optimal-smooth}.

\begin{lemma}[Bounded trigonometric approximation]
\label{lem:classical-trigonometric-density}
Let $m^*=\pi(z)$, where $z$ is a finite descendant formula with
Gaussian leaf seeds.  Assume that every transmitted child column
is odd, and that each local map producing such a column or the
output is a jointly odd smooth function of its local seed and
finite list of odd aggregate inputs, with polynomial-growth
derivatives.  Purely local intermediate operations may be composed
into these maps before the lemma is applied.
For every $\eta>0$, there is an odd bounded angle formula $\theta$
such that
\begin{equation}
 \norm{\sin(2\theta)-m^*}<\eta.
 \label{eq:angle-density}
\end{equation}
Every sent column in the formula for $\theta$ can be chosen bounded
and odd.  Its pointwise maps can be chosen to be finite real
trigonometric sums in their finite lists of odd inputs, with odd
joint parity for odd outputs.  Its local seed inputs can be chosen
bounded and odd.  All these functions have bounded derivatives of
every order in their continuous arguments.
\end{lemma}

\begin{proof}
Replace a Gaussian local seed $s$ by a smooth bounded odd truncation
$s_R$ satisfying $|s_R|\le|s|$ and $s_R\to s$ pointwise.  It converges
in every finite moment.  At a fixed finite depth this changes the
original formula by an amount tending to zero in every required
finite moment: use \eqref{eq:classical-gaussian-isometry} at each
aggregation, Gaussian moment bounds for its finite list of fields,
and the polynomial bounds for the derivatives of the pointwise maps.

Fix $0<\kappa<1$.  The desired final angle
\[
 \theta_\kappa=\tfrac12\arcsin((1-\kappa)m^*)
\]
is bounded and odd, and is a Lipschitz function of the clipped output.
It satisfies $\sin(2\theta_\kappa)=(1-\kappa)m^*$.
We approximate this angle and all its subformulas beginning at the
output and working toward the leaves.  For the finite-dimensional
ideal input law at a pointwise operation, a continuous polynomial-growth
function has an arbitrarily accurate $L^p$ approximation, for any
fixed finite $p$, by a bounded smooth function with bounded
derivatives.  This follows by a smooth cutoff outside a large box
and mollification; the uniform moments control the cutoff error.
Symmetrizing by $F(x)\mapsto(F(x)-F(-x))/2$ preserves oddness and
does not increase its error under the symmetric input law.

Once an outer map has been replaced by a fixed bounded Lipschitz
map, perturbations of its inputs propagate continuously in $L^2$.
For its aggregate inputs this is exactly
\eqref{eq:classical-gaussian-isometry}, jointly for the finite list
of child columns.  We can therefore choose the approximations at
the next lower level sufficiently accurately to preserve the
already chosen output accuracy.  There are finitely many operations,
so induction supplies a formula in which every transmitted scalar
function is bounded, smooth, and odd.  No uniform bound on Lipschitz
constants over the sequence of approximation choices is needed;
each outer choice is fixed before its inputs are approximated.

Finally a bounded smooth cutoff supported strictly inside a box can
be extended periodically to a larger box.  Fourier approximation of
this smooth periodic function gives finite trigonometric sums
converging uniformly, with any prescribed finite number of
derivatives, on that box.  To control the complement, first fix the
box so that its probability is small and choose the Fourier
approximation uniformly on the entire period; both the cutoff and
its approximant then have a common finite sup norm.  The same
output-to-leaf continuity argument applies.  Odd symmetrization
leaves only sine terms $\sin(\lambda\cdot x)$ for an odd scalar
output.  This constructs the asserted trigonometric formula.
Choose its angle error smaller than $\eta/4$ and $\kappa<\eta/2$.
Since $x\mapsto\sin(2x)$ is $2$-Lipschitz, the resulting output
satisfies \eqref{eq:angle-density}.
\end{proof}

\subsection{The synthesis interface}

A \emph{mask} is an event in the local seed space invariant under seed
sign reversal.  Masks at different vertices are evaluated on their
respective seeds.  We use a finite partition into masks and write
\(p_R\) for the probability of a mask \(R\).  Probabilities used as
denominators are positive.  For disjoint masks \(R,S\), let
\[
 (A_{RS}f)_i=\1_R(i)D^{-1/2}
       \sum_{j\sim i}\1_S(j)f_j .
\]
This notation first denotes a finite-tree operation.  Its rooted limit
is provided by Theorem~\ref{thm:gate-calculus}.

The scalar formulas considered here have finite expression trees.
Their leaves are bounded measurable functions of the local seed.
An internal operation is either a pointwise function of finitely many
already constructed quantities or an aggregation \(A_{RS}\).
Pointwise functions are smooth in their aggregate arguments, with
polynomial bounds on every derivative, uniformly in the local seed.
Local seed dependence itself need not be smooth.  All quantities sent
across edges are odd under simultaneous reversal of seeds and odd spin
axes.  In the classical formulas there are no spin arguments, so this
is simply joint seed parity.  Sent nonlinear functions, and the output
angles, are bounded.  The aggregate fields need not be bounded.

\begin{definition}[Exclusion formulas]
\label{def:exclusion-formulas}
An aggregation from \(S\) to \(R\) is \emph{excluded} if the formula
sent from \(S\) uses no vertex in \(R\) anywhere below that
aggregation.  The formulas used in this section satisfy this
requirement at every aggregation and retain all ancestor exclusions:
the masks along any branch of an expression tree are distinct.
Different occurrences of a formula with different exclusion lists are
regarded as different formulas.  A gate target is excluded from all
cross-edge dependencies of its parameters.
\end{definition}

The distinction between an expression tree and the underlying regular
tree is useful.  A formula may reuse a mask in different expression
branches.  The exclusion rule forbids its reuse along a single active
branch.  It therefore forbids an immediate return across a cut, without
requiring independent values for two formulas that happen to use the
same descendant subtree.

\subsection{Colors, exclusions, and helpers}

A descendant formula omits the parent \emph{vertex} in its finite-tree
interpretation. We now impose the stronger type exclusions needed
for synthesis. We estimate their cost in the limiting classical space,
where \eqref{eq:classical-gaussian-isometry} makes the cost of deleting
small-probability types explicit.

Let $s_{\mathrm{raw}}$ denote the single standard Gaussian available
as the local control seed. Its sign is an independent fair sign, and
the distribution function of its absolute value is uniform on $(0,1)$.
Splitting the
binary digits of that uniform variable into disjoint subsequences
produces independent uniform variables, apart from a null set of
dyadic ambiguities.  Use one for a finite color variable $C$, and
another, together with the original sign, to generate a fresh
standard Gaussian $S$. This defines a measurable map
\begin{equation}
 \begin{gathered}
 s_{\mathrm{raw}}\longmapsto
       \bigl(C(s_{\mathrm{raw}}),S(s_{\mathrm{raw}})\bigr),\\
 C(-s_{\mathrm{raw}})=C(s_{\mathrm{raw}}),\qquad
 S(-s_{\mathrm{raw}})=-S(s_{\mathrm{raw}}),
 \end{gathered}
 \label{eq:raw-seed-reencoding}
\end{equation}
whose image has the product law of the chosen color and a standard
Gaussian. Every colored formula below is evaluated at this pair at
each vertex: its Gaussian input is $S(s_{\mathrm{raw}})$.
The necessary color masks are bounded measurable even functions of
the original local seed.  Smoothness below refers to the continuous
message arguments conditional on a color; no smoothness of a mask
as a function of the original seed is asserted or required.

One may first adjoin an independent color to the color-free seed
space and then pull back by \eqref{eq:raw-seed-reencoding}.
The local pullback $f(C,S)\mapsto
f(C(s_{\mathrm{raw}}),S(s_{\mathrm{raw}}))$ preserves the product
law, inner products, and the constant seed vector. Tensoring it at
the root and applying the symmetric Fock functor to its centered
child restriction extends it to every recursive height. These
isometries respect the empty-descendant embeddings, preserve parity,
and intertwine $\cL$, which inserts a constant seed at the new root.
Thus both adjoining colors and pulling the entire formula back to
the raw seeds preserve its rooted energy. Partition the colors
into nonhelper types $1,\ldots,M$ and helper types.  Write
\[
 \max_{1\le c\le M}\Pr(C=c)\le\alpha,
 \qquad \Pr(C\text{ is a helper})=p_H.
\]
Every helper type may have a prescribed positive probability, subject
to the total being sufficiently small.  The number of helper types
needed later is finite, at most one per nested propagation depth.

\begin{lemma}[Type exclusion at small cost]
\label{lem:classical-color-exclusion}
Fix a descendant formula in which every transmitted column and the
output are given by jointly odd bounded smooth functions of the
local odd seed and the odd aggregate inputs, with bounded first
derivatives.  Let its
propagation depth be $h$.  It has a finite realization by colored
formulas such that every aggregation excludes the current color,
all colors of the path ancestors, and all helper colors.
If $f$ is its original output and $f^{\mathrm{col}}$ its colored
output, set to zero on helper roots, then
\begin{equation}
 \norm{f^{\mathrm{col}}-f}^2\le C(\alpha+p_H),
 \label{eq:color-exclusion-bound}
\end{equation}
where $C$ depends only on the fixed color-free formula, not on the
number or individual probabilities of its nonhelper colors.
Conditional bounds for the unchanged nonlinear pointwise functions
and their derivatives in the original grouped aggregate coordinates
are uniform over these colors.  No bound independent of $M$ is
asserted for the Euclidean derivative norm after flattening all
source-color fields into separate coordinates.
\end{lemma}

\begin{proof}
Unroll the fixed expression into a finite expression tree.  For
each occurrence, each possible current nonhelper color $c$, and each
possible set $A$ of ancestor colors, make a copy of the lower
formula.  Its aggregates retain only children whose colors avoid
$A\cup\{c\}$ and all helpers.  At a retained child of color $d$,
the recursive copy uses ancestor set $A\cup\{c\}$ and current color
$d$.  Thus its own aggregates exclude $A\cup\{c,d\}$.  All sets
encountered have size at most $h+1$.  There are finitely many copies
because the original depth and the color partition are finite.
The aggregation at color $c$ is the sum of its source-color
aggregations.  Linearity of $\mathsf G$ identifies this finite sum
with the lift of the corresponding masked child column.

These copies also define an orientation-free all-neighbor rule.
When a recursive call crosses from a vertex $i$ of color $c$ to a
vertex $j$ of color $d$, the ancestor list at $j$ contains $c$.
Every aggregate at $j$ therefore gives the neighbor $i$ a zero
mask. Its sum over all neighbors equals its sum over neighbors other
than $i$, with the same normalization $D^{-1/2}$. Induction over the
expression depth identifies the all-neighbor calculation with the
descendant calculation on each active branch. The outer root has no
parent and uses all its neighbors. In particular the finite rule uses
the aggregations $A_{RS}$ of the synthesis interface without requiring
a parent orientation. A full root has $D+1$ neighbors; its limiting
formula agrees with the $D$-child convention by the full-root assertion
of Theorem~\ref{thm:gate-calculus}.

Here is a uniform error estimate.  Give the original formula the
extra independent colors but ignore them in its computation.
At one aggregate, let $g$ be its ideal bounded child column, with
$|g|\le K$, and let $g^{A,d}$ be the recursively modified column on
a retained color $d$.  Couple the two fields using
\eqref{eq:classical-gaussian-isometry}.  The modified aggregate has
child column
\[
 \widetilde g
 =\sum_{d\notin A\cup\{c\},\ d\ {\mathrm{nonhelper}}}
            \1_{\{C=d\}}g^{A\cup\{c\},d}.
\]
The color masks are orthogonal, and $g$ is independent of the colors
at its root.  Hence
\begin{align*}
 \norm{\widetilde g-g}^2
 &\le K^2\bigl(p_H+(|A|+1)\alpha\bigr)\\
 &\quad+\max_{d\notin A\cup\{c\}}
       \E\bigl[|g^{A\cup\{c\},d}-g|^2\mid C=d\bigr].
\end{align*}
Every masked column is centered: color masks are even and the
conditional formula is odd in the Gaussian seeds.  Thus the
Gaussian-isometry identity applies without a mean correction.
For a list of $k$ aggregates the squared Euclidean input error is
the sum of their squared errors.  A pointwise map with Lipschitz
constant $K'$ multiplies that sum by at most $(K')^2$.
Inducting over the finite expression depth, with the unchanged
bounded local seed as the base case, proves a conditional squared
error bounded by $C(\alpha+p_H)$ uniformly over current and
ancestor colors.  Setting the helper-root output to zero adds at
most $\norm f_\infty^2p_H$.

This estimate has the fan-in of the original formula, not the number
of color copies: disjoint child masks were combined by squared
norms before applying a Lipschitz bound.  Each nonlinear pointwise
map still receives the original finite list of grouped fields: if
it originally had $k$ field arguments, its colored version is
$F(S,Y_1^{\mathrm{col}},\ldots,Y_k^{\mathrm{col}})$, where
\[
 Y_r^{\mathrm{col}}=\sum_d\mathsf G(\1_{\{C=d\}}g_{r,d})
       =\mathsf G\left(\sum_d\1_{\{C=d\}}g_{r,d}\right).
\]
The conditional derivative bounds in these $k$ arguments are
unchanged.  The finite sums defining grouped fields are additional
linear operations and can be unbounded; boundedness is required of
the transmitted columns and final angles, not of Gaussian fields.
For a trigonometric formula, substitute each grouped linear sum into
the next pointwise map.  A trigonometric term in linear combinations
of the source fields is still a trigonometric term, so this substitution
preserves the finite trigonometric form without retaining separate
linear pointwise maps.
For the synthesis interface, a rotation by a grouped field is
exactly the product of the source-field rotations on the same axis,
since these classical fields commute.  If a compilation instead
regards all the source fields as separate coordinates, its
derivative constants may depend on the fixed number $M$.  This is
harmless: $M$ and all these finite lists are fixed before any helper
refinement or quantum approximation is taken.  Only the classical
error constant in \eqref{eq:color-exclusion-bound} and the bounds in
the original grouped coordinates are asserted to be independent
of $M$.
\end{proof}

\begin{proposition}[Classical near-optimality in synthesis form]
\label{prop:classical-near-optimum}
For every $\varepsilon>0$ and every $\eta_H>0$, there are a finite
color partition with helper probability $0<p_H<\eta_H$, a finite
number of helper types, and finite bounded odd angle formulas
$\theta_c$ on the nonhelper colors such that the following hold.
Each transmitted column is bounded and odd, every pointwise map
may be chosen as a finite trigonometric sum in its odd inputs, and
every aggregation excludes its current color, all path-ancestor
colors, and all helpers.  Conditional on the finite labels, all
continuous derivatives of the pointwise maps are bounded.  The
output
\begin{equation}
 m^{\mathrm{col}}
  =\sum_{c=1}^M\1_{\{C=c\}}\sin(2\theta_c)
 \label{eq:final-classical-formula}
\end{equation}
is centered, satisfies $|m^{\mathrm{col}}|\le1$, vanishes on helper
roots, and obeys
\begin{equation}
 \cE(m^{\mathrm{col}})\ge P_*-\varepsilon.
 \label{eq:final-classical-value}
\end{equation}
After the uncolored finite formula has been fixed, the total helper
probability may be taken arbitrarily small, with all individual
helper probabilities positive.  All depths, coefficients, label laws,
and functions in the conclusion are fixed independently of degree,
system size, and disorder.
\end{proposition}

\begin{proof}
Use \eqref{eq:scalar-near-optimal-smooth} with error
$\varepsilon/3$ and write $m^*=z$.  Apply Lemma~\ref{lem:classical-trigonometric-density}
to obtain an angle $\theta$ whose output $m=\sin(2\theta)$ has
$\norm{m-m^*}<\varepsilon/12$.  Both outputs are bounded by one,
so \eqref{eq:energy-l2-continuity} loses less than
$\varepsilon/6$.  Fix this entire bounded trigonometric formula,
including its finite depth and Lipschitz constants.

Reserve a distinct helper type for each propagation depth.
Choose their positive probabilities with total
$p_H<\eta_H$ sufficiently small, and partition the remaining
probability into finitely many colors of maximum probability
$\alpha$ sufficiently small.  Apply
Lemma~\ref{lem:classical-color-exclusion} to the angle calculation
and use the $2$-Lipschitz property of sine.  Taking $\alpha,p_H$
smaller if necessary gives
$\norm{m^{\mathrm{col}}-m}<\varepsilon/6$.
Another application of \eqref{eq:energy-l2-continuity} loses less
than $\varepsilon/3$.  The total loss is less than $\varepsilon$.
Oddness gives centering, and the remaining claims follow directly
from the construction. Every choice has been made for a fixed finite
calculation in the limiting rooted space; no uniform rate in
the number of colors, circuit depth, or helper probabilities has
been used.
\end{proof}
